\input{../../resources/lesson-head.tex}

\subsection*{Week 8 at a glance}

\subsubsection*{We will be learning and practicing to:}
%classify cardinality
%important sets
%function and relation definitions
%function properties
%functions for cardinality
%contradiction proofs
\begin{itemize}

\item Clearly and unambiguously communicate computational ideas using appropriate formalism. Translate across levels of abstraction.
\begin{itemize}
    \item Defining important sets of numbers, e.g. set of integers, set of rational numbers
    \item Classifying sets into: finite sets, countably infinite sets, uncountable sets
    \item Defining functions, predicates, and binary relations using multiple representations
    \item Determining whether a given binary relation is symmetric, antisymmetric, reflexive, and/or transitive
    \item Determining whether a given binary relation is an equivalence relation and/or a partial order
\end{itemize}

\item Know, select and apply appropriate computing knowledge and problem-solving techniques. Reason about computation and systems. Use mathematical techniques to solve problems. Determine appropriate conceptual tools to apply to new situations. Know when tools do not apply and try different approaches. Critically analyze and evaluate candidate solutions.
\begin{itemize}
    \item Determining what evidence is required to establish that a quantified statement is true or false
    \item Evaluating quantified statements about finite and infinite domains
\end{itemize}

\item Apply proof strategies, including direct proofs and proofs by contradiction, and determine whether a proposed argument is valid or not.
\begin{itemize}
    \item Using proofs as knowledge discovery tools to decide whether a statement is true or false
\end{itemize}
\end{itemize}

\subsubsection*{TODO:}
\begin{list}
   {\itemsep2pt}
    \item Review quiz based on Week 7 class material (due Monday)
    \item Homework 4 due Tuesday
    \item Test 2 in the CBTF at your scheduled time!
You will be required to show a physical photo identification in order
to test. Electronic IDs and photos of IDs are not accepted. Additionally, you will only be allowed
to bring your ID and a writing utensil into the testing room. You will have to store away all of
your other belongings before entering the room. Please see the Triton Testing Center's website
for more details.
\end{list}

\newpage


\section*{Week 8 Monday and Wednesday: Finite, countably infinite, and uncountable sets}
\subsection*{Cardinality of sets}
\input{../activity-snippets/cardinality-definition.tex}
\input{../activity-snippets/cardinality-caution.tex}
%! app: TODOapp
%! outcome: functions for cardinality

{\bf Properties of cardinality}
\begin{align*}
&\forall A ~ (~  |A| = |A| ~)\\
&\forall A ~ \forall B ~(~ |A| = |B|  ~\to ~ |B| = |A|~)\\
&\forall A ~ \forall B ~ \forall C~ (~ (|A| = |B| ~\wedge~ |B| = |C|) ~\to ~ |A| = |C|~)
\end{align*}

{\it Extra practice with proofs:} Use the definitions of bijections to prove these properties.
\input{../activity-snippets/cantor-schroder-bernstein-theorem.tex}
\newpage
\input{../activity-snippets/countably-infinite-definition.tex}
\input{../activity-snippets/countably-infinite-examples-sets-of-numbers.tex}
%! app: TODOapp
%! outcome: TODOoutcome

{\bf More examples of countably infinite sets}

{\bf Claim}: $S$ is countably infinite

{\it Similarly: The set of all strings over a specific alphabet is countably infinite.}
\begin{center}
Bijection using alphabetical-ish ordering
(first order by length, then alphabetically among strings of same length) 
of strands
\end{center}

\vspace{50pt}


{\bf Claim}: $L$ is countably infinite

\begin{multicols}{2}
\begin{align*}
    &list: \mathbb{N} \to L \\
    &list(n) = (n, []) \\
    &
\end{align*}

\begin{align*}
    &toNum: L \to \mathbb{N} \\
    &toNum([]) = 0 \\
    &toNum( ~(n,l)~) = 2^n 3^{toNum(l)} \qquad \text{for $n \in \mathbb{N}$, $l \in L$}
\end{align*}
\end{multicols}
\vspace{30pt}

{\bf Claim}: $|\mathbb{Z}^+| = |\mathbb{Q}|$ 

One-to-one function from $\mathbb{Z}^+$ to $\mathbb{Q}$
is $f_1: \mathbb{Z}^+ \to \mathbb{Q}$ with $f_1(n) = n$ for all 
$n \in \mathbb{N}$.

\vspace{30pt}


\begin{align*}
    &f_2: \mathbb{Q} \to \mathbb{Z} \times \mathbb{Z} \\
    &f_2(x) = \begin{cases}
        (0,1) & \text{if $x = 0$} \\
        (p,q) & \text{if $x = \frac{p}{q}$,}\\
                & \text{$gcd(p,q) = 1$, $q > 0$}
    \end{cases}
\end{align*}
\begin{multicols}{2}
\begin{align*}
    &f_3: \mathbb{Z} \times \mathbb{Z} \to \mathbb{Z}^+ \times \mathbb{Z}^+ \\
    &f_3(~(x,y)~) = \begin{cases}
        (2x+2, 2y+2) & \text{if $x \geq 0, y \geq 0$} \\
        (-2x-1, 2y+2) & \text{if $x < 0, y \geq 0$} \\
        (2x+2, -2y+1) & \text{if $x \geq 0, y < 0$} \\
        (-2x-1, -2y-1) & \text{if $x < 0, y < 0$} \\
    \end{cases}
\end{align*}

\begin{align*}
    &f_4: \mathbb{Z}^+ \times \mathbb{Z}^+ \to \mathbb{Z}^+ \\
    &f_4(~(x,y)~) = 2^x 3^y \qquad \text{for $x,y \in \mathbb{Z}^+$}
\end{align*}
\end{multicols}
\newpage
\subsection*{Cardinality categories}
\input{../activity-snippets/cardinality-categories.tex}
\input{../activity-snippets/cardinality-countability-lemmas.tex}

\subsection*{Are there always *bigger* sets?}
%! app: TODOapp
%! outcome: functions for cardinality, classify cardinality

{\it Recall}: When $U$ is a set, $\mathcal{P}(U) = \{ X \mid X \subseteq U\}$

{\it Key idea}: For finite sets, the power set of a set has strictly greater size than the set itself.
Does this extend to infinite sets?

{\bf Definition}: For two sets $A, B$, we use the notation $|A| < |B|$ to denote
$(~|A| \leq |B| ~) \land \lnot (~|A| = |B|)$.

\begin{alignat*}{4}
    &\emptyset = \{ \} \qquad &&\mathcal{P}(\emptyset) = \{ \emptyset \} \qquad &&|\emptyset| < |\mathcal{P}(\emptyset)| \\
    &\{1 \} \qquad &&\mathcal{P}(\{1\}) = \{ \emptyset, \{1\} \} \qquad &&|\{1\}| < |\mathcal{P}(\{1\})| \\
    &\{1,2 \} \qquad &&\mathcal{P}(\{1,2\}) = \{ \emptyset, \{1\}, \{2\}, \{1,2\} \} \qquad &&|\{1,2\}| < |\mathcal{P}(\{1,2\})| \\
\end{alignat*}

{\bf $\mathbb{N}$ and its power set}

Example elements of $\mathbb{N}$ 

\vspace{20pt}

Example elements of $\mathcal{P}(\mathbb{N})$

\vspace{20pt}

{\bf Claim}: $| \mathbb{N} | \leq |\mathcal{P} ( \mathbb{N} ) |$

\vspace{100pt}
\newpage
{\bf Claim}: There is an uncountable set.  Example: $\underline{\phantom{~~~\mathcal{P}(\mathbb{N})~~~}}$

{\bf Proof}:  By definition of countable, since $\underline{\phantom{~~~\mathcal{P}(\mathbb{N})~~~}}$
is not finite, {\bf to show} is $|\mathbb{N}| \neq  |\mathcal{P}(\mathbb{N})|$ .

Rewriting using  the definition of  cardinality, {\bf to show} is

\phantom{$\neg \exists f : \mathbb{N} \to \mathcal{P}(\mathbb{N})  ~~(f \text{ is a bijection})~~$}

\phantom{or equivalently $\forall f : \mathbb{N} \to \mathcal{P}(\mathbb{N})  ~~(f \text{ is not a bijection})~~$}


Towards a proof by  universal generalization,  consider  an arbitrary function $f:  \mathbb{N} \to\mathcal{P}(\mathbb{N})$.

{\bf To show}: $f$ is not a bijection.  It's enough to show that $f$ is not onto.

Rewriting using the definition of  onto, {\bf to show}:
\[
\neg  \forall  B \in  \mathcal{P}(\mathbb{N}) ~\exists a \in \mathbb{N}  ~(~f(a) =  B~)
\]
. By logical  equivalence, we can write this as an existential statement:
\[
\underline{\phantom{\qquad\qquad\exists B \in  \mathcal{P}(\mathbb{N}) ~\forall a \in \mathbb{N}  ~(~f(a) \neq  B~)\qquad\qquad}}
\]
In search of a witness, define the following  collection of nonnegative integers:
\[
D_f = \{ n \in \mathbb{N}  ~\mid~  n \notin f(n)  \}
\]
. By  definition  of power  set, since  all elements  of  $D_f$ are  in  $\mathbb{N}$,   $D_f \in \mathcal{P}(\mathbb{N})$.  It's enough to prove the following Lemma: 

{\bf Lemma}: $\forall a \in \mathbb{N}  ~(~f(a) \neq  D_f~)$.


{\bf Proof  of lemma}: \phantom{Towards universal  generalization, consider an arbitrary  $a \in \mathbb{N}$.
By definition  of set equality, {\bf to show} is  $\exists  x ( \neg  (x \in f(a)~  \leftrightarrow  ~x \in D_f))$.
For a witness, consider $x = a$.  There are two cases:  $a \in  f(a)~\vee~a \notin f(a)$. By definition 
of $D_f$, each guarantees that $f(a) \neq  D_f$.}\\

\vspace{50pt}

By  the Lemma, we  have proved that $f$ is not onto, and since $f$ was arbitrary, there are no onto
functions from $\mathbb{N}$ to $\mathcal{P}(\mathbb{N})$. QED


{\bf Where does $D_f$ come from?} The idea is to build a set that would ``disagree" with 
each of the images of $f$ about some element. 

\begin{center}
\begin{tabular}{c|c|ccccccc}
$n \in \mathbb{N}$ & $f(n) = X_n$ &  Is $0   \in X_n$?   & Is $1 \in X_n$?  &  Is $2 \in X_n$?  &  Is $3 \in X_n$?  &
 Is $4 \in X_n$?  &  \ldots & Is $n \in D_f$?\\
\hline
$0$ & $f(0) = X_0$ & {\bf  Y~/~N}  & Y~/~N & Y~/~N & Y~/~N &Y~/~N & \ldots & {\bf  N~/~Y }\\
$1$ & $f(1) = X_1$ & Y~/~N  & {\bf  Y~/~N} & Y~/~N & Y~/~N & Y~/~N & \ldots & {\bf  N~/~Y }\\
$2$ & $f(2) = X_2$ & Y~/~N  & Y~/~N & {\bf  Y~/~N} & Y~/~N &Y~/~N & \ldots & {\bf  N~/~Y }\\
$3$ & $f(3) = X_3$ & Y~/~N  & Y~/~N & Y~/~N & {\bf  Y~/~N} & Y~/~N & \ldots & {\bf  N~/~Y }\\
$4$ & $f(4) = X_4$ & Y~/~N  & Y~/~N & Y~/~N & Y~/~N &{\bf  Y~/~N} & \ldots & {\bf  N~/~Y }\\
\vdots
\end{tabular}
\end{center}
\newpage
\subsection*{Countable vs.\ uncountable: sets of numbers}
%! app: TODOapp
%! outcome: functions for cardinality, classify cardinality, important sets

{\bf Comparing $\mathbb{Q}$ and $\mathbb{R}$} 


Both $\mathbb{Q}$ and $\mathbb{R}$ have no greatest element.

Both $\mathbb{Q}$ and $\mathbb{R}$ have no least element.

The quantified statement 
\[
    \forall x \forall y (x < y \to \exists z ( x < z < y) )
\]
is true about both $\mathbb{Q}$ and $\mathbb{R}$.

Both $\mathbb{Q}$ and $\mathbb{R}$ are infinite. But, $\mathbb{Q}$ is countably infinite
whereas $\mathbb{R}$ is uncountable.\\


{\bf The set of real numbers}

$\mathbb{Z} \subsetneq \mathbb{Q} \subsetneq \mathbb{R}$


{\bf  Order axioms} (Rosen Appendix 1): 

\begin{center}
\begin{tabular}{p{1.2in}p{4in}}
Reflexivity &  $\forall a \in  \mathbb{R} (a \leq a)$\\
Antisymmetry &  $\forall a \in  \mathbb{R}~\forall b \in \mathbb{R}~(~(a \leq b~ \wedge ~b \leq a) \to (a=b)~)$\\
Transitivity &  $\forall a \in  \mathbb{R}~\forall b \in \mathbb{R}~\forall c \in \mathbb{R}~
(~(a \leq b \wedge b \leq c) ~\to  ~(a \leq c)~)$ \\
Trichotomy & 
$\forall a \in \mathbb{R}~\forall b \in \mathbb{R}~ ( ~(a=b ~\vee~ b > a ~\vee~ a  < b)  $
\end{tabular}
\end{center}


{\bf  Completeness axioms} (Rosen Appendix 1): 


\begin{center}
\begin{tabular}{p{1.4in}p{6in}}
Least upper bound &  Every nonempty set of real numbers that 
is bounded  above has  a  least upper bound  
%\phantom{~$\forall X \in \mathcal{P}(\mathbb{R}) ( ~(\exists x (x \in X)  \wedge \exists  y ~\forall x ( x \in X ~\to~ x \leq y) )~\to ~\exists  y_0 \left( ~\forall x ( x \in X ~\to~ x \leq y) \wedge 
%\forall z ( \forall  x (x \in X ~\to~ x \leq z) ~\to ~ y  \leq  z)\right) )$}\\
\\
Nested intervals &  For each sequence  of intervals  $[a_n , b_n]$
where, for each $n$, $a_n < a_{n+1} < b_{n+1} < b_n$, there
is at least one  real number $x$ such that, for all $n$, 
$a_n \leq x \leq b_n$.\\
\end{tabular}
\end{center}

Each real  number $r  \in  \mathbb{R}$ is described by a function to give better and better approximations
\[
x_r: \mathbb{Z}^+ \to \{0,1\}  \qquad  \text{where  $x_r(n ) =  n^{th} $ bit in  binary expansion of $r$}
\]
\begin{center}
\begin{tabular}{|c|c|p{3.9in}|}
\hline
$r$ & Binary expansion & $x_r$ \\
\hline
%$0.5$ & $0.10000\ldots$ &  $x_{0.5}(n) = \begin{cases} 1 &\text{if $n = 1$} \\ 0&\text{otherwise} \end{cases}$\\
%& & \\
%\hline
$0.1$ & $0.00011001 \ldots$ &  $x_{0.1}(n) = \begin{cases} 0&\text{if $n > 1$ and $(n~\text{\bf mod}~4) =2$} \\
0&\text{if $n=1$ or if $n > 1$ and $(n~\text{\bf mod}~4) =3$} \\1&\text{if $n > 1$ and $(n~\text{\bf mod}~4) =0$} \\
1&\text{if $n > 1$ and $(n~\text{\bf mod}~4) =1$} \end{cases}$  \\
&&  \\
\hline
$\sqrt{2} - 1 = 0.4142135 \ldots$  &$0.01101010\ldots$& Use linear approximations
(tangent lines from calculus) to get algorithm for bounding error of successive operations. Define 
$x_{\sqrt{2}-1}(n)$ to be  $n^{th}$ bit in approximation  that has error less than  $2^{-(n+1)}$.
\\
&& \\
\hline
\end{tabular}
\end{center}

\newpage 

{\bf Claim}: $\{  r \in \mathbb{R} ~\mid~ 0 \leq r ~\wedge~ r \leq 1 \}$ is uncountable.

{\it Approach 1}: Mimic proof that $\mathcal{P}(\mathbb{Z}^+)$ is uncountable.


{\bf Proof}:  By definition of countable, since $\{  r \in \mathbb{R} ~\mid~ 0 \leq r ~\wedge~ r \leq 1 \}$
is not finite, {\bf to show} is $|\mathbb{N}| \neq  |\{  r \in \mathbb{R} ~\mid~ 0 \leq r ~\wedge~ r \leq 1 \}|$ .


{\bf To show} is
$\forall f : \mathbb{Z}^+ \to \{  r \in \mathbb{R} ~\mid~ 0 \leq r ~\wedge~ r \leq 1 \}  ~~(f \text{ is not a bijection})~~$.
Towards a proof by  universal generalization, consider  an arbitrary function 
$f:  \mathbb{Z}^+ \to \{  r \in \mathbb{R} ~\mid~ 0 \leq r ~\wedge~ r \leq 1 \}$.
{\bf To show}: $f$ is not a bijection.  It's enough to show that $f$ is not onto.
Rewriting using the definition of  onto, {\bf to show}:
\[
\exists x \in \{  r \in \mathbb{R} ~\mid~ 0 \leq r ~\wedge~ r \leq 1 \} ~\forall a \in \mathbb{N}  ~(~f(a) \neq  x~)
\]
In search of a witness, define the following  real number by defining its binary expansion
\[
d_f = 0.b_1 b_2 b_3 \cdots
\]
where $b_i = 1-b_{ii}$ where $b_{jk}$ is the coefficient of $2^{-k}$ in the binary expansion of $f(j)$.
Since\footnote{There's a subtle imprecision in this part of the proof as presented, but it can be fixed.} $d_f \neq f(a)$ for any positive integer $a$, $f$ is not onto.
%$x_r: \mathbb{Z}^+ \to  \{0,1\}$ where  $x_r(n) = n^{th}$ bit in binary expansion  of $r$
%\qquad 
%$X_r =  \{ n \in \mathbb{Z}^+ ~\mid~ x_r(n) = 1 \}$

{\it Approach 2}: Nested closed interval property

{\bf To show} $f: \mathbb{N} \to \{  r \in \mathbb{R} ~\mid~ 0  \leq r ~\wedge~ r \leq 1 \}$ is not onto. 
{\bf  Strategy}: Build
a sequence of nested closed intervals that each avoid some $f(n)$.   Then  the real
number that is in all of the intervals  can't be $f(n)$ for any $n$. Hence,  $f$ is not  onto.

Consider the function $f: \mathbb{N} \to \{  r \in \mathbb{R} ~\mid~ 0 \leq r ~\wedge~ r \leq 1 \}$ with  $f(n) = \frac{1+\sin(n)}{2}$

\begin{center}
\begin{tabular}{c||p{1.65in} || p{3in} }
$n$ &  $f(n)$& Interval that avoids $f(n)$ \\
\hline
$0$ & $0.5$ &  \\
$1$ &$0.920735\ldots$  &  \\
$2$ &$0.954649\ldots$ &  \\
$3$ &$0.570560\ldots$ & \\
$4$ &$0.121599\ldots $&  \\
\vdots &  &\\
\end{tabular}
\end{center}
 
\subsection*{Other examples of uncountable sets}
%! app: TODOapp
%! outcome: classify cardinality

\begin{itemize}
    \item The power set of any countably infinite set is uncountable. For example:
    \[
        \mathcal{P}(\mathbb{N}), \mathcal{P}(\mathbb{Z^+}), \mathcal{P}(\mathbb{Z})
    \]
    are each uncountable.
    \item The closed interval $\{x \in \mathbb{R} ~|~ 0 \leq x \leq 1\}$, any other nonempty closed interval of real numbers whose endpoints are 
    unequal, as well as the related intervals that exclude one or both of the endpoints.
    \item The set of all real numbers $\mathbb{R}$ is uncountable and the set of irrational
    real numbers $\overline{\mathbb{Q}}$ is uncountable.
\end{itemize}

\newpage
\section*{Week 8 Friday: Cardlinality and binary relations}

\subsection*{Cardinality of sets: recap}
%! app: TODOapp
%! outcome: functions for cardinality, classify cardinality, important sets


The set of positive integers $\mathbb{Z}^{+}$ is countably infinite.

The set of integers $\mathbb{Z}$ is countably infinite and is a proper superset of $\mathbb{Z}^{+}$. 
In fact, the set difference 
$$\mathbb{Z} \setminus \mathbb{Z}^{+} = \{ x \in \mathbb{Z} \mid x \notin \mathbb{Z}^+\} 
= \{ x \in \mathbb{Z} \mid x \leq 0 \}$$ is countably infinite.

The set of rationals 
$\mathbb{Q} = \left\{ \frac{p}{q} \mid p \in \mathbb{Z}  \text{ and  } 
q  \in \mathbb{Z} \text{ and } q \neq  0 \right\}$ is countably infinite.


The set of real numbers $\mathbb{R}$ is uncountable. In fact, the closed 
interval $\{x \in \mathbb{R} ~|~ 0 \leq x \leq 1\}$, 
any other nonempty closed interval of real numbers whose endpoints are 
unequal, as well as the related intervals that exclude one or both of the endpoints are
each uncountable.
The set of {\bf irrational} numbers $\overline{\mathbb{Q}} = \mathbb{R} - \mathbb{Q}
= \{ x \in \mathbb{R} \mid x \notin \mathbb{Q} \}$ is uncountable.

\vfill

We can classify any set as 
\begin{itemize}
\item {\bf Finite} size. Fact: For each positive number $n$, for any sets $X$ and $Y$ each 
size $n$, there is a bijection between $X$ and $Y$.
\item {\bf Countably Infinite}. Fact: for any countably infinite sets $X$ and $Y$, there is a bijection
between $X$ and $Y$.
\item {\bf Uncountable}. Examples: $\mathcal{P}(\mathbb{N})$, the power set of any infinite set, 
the set of real numbers, any nonempty interval of real numbers. Fact: there are (many) examples 
of uncountable sets that do not have a bijection between them.
\end{itemize}
\newpage
%! app: TODOapp
%! outcome: TODOoutcome

{\bf True or False}?  For all sets $A$ and $B$ where  $A \subseteq B$, 
if $A$ is infinite then $B$ is finite.

\vfill

{\bf True or False}?   For all sets $A$ and $B$ where  $A \subseteq B$, 
if $A$ is countable then $B$ is countable.

\vfill

{\bf True or False}?   For all sets $A$ and $B$ where  $A \subseteq B$, 
if $B$ is infinite then $A$ is finite.

\vfill

{\bf True or False}?   For all sets $A$ and $B$ where  $A \subseteq B$, 
if $B$ is uncountable then $A$ is countable.

\vfill
\newpage

\newpage
\subsection*{Binary relations}
\input{../activity-snippets/binary-relation-definition.tex}
\vfill
%! app: TODOapp
%! outcome: graph representations of relations

For relation $R$ on a set $A$, we can represent this relation as a
{\bf graph}: a collection of nodes (vertices) and edges (arrows). The 
nodes of the graph are the elements of $A$ and 
there is an edge from $a$ to $b$ exactly when $(a,b) \in R$.

\vfill
\newpage
%! app: TODOapp
%! outcome: TODOoutcome

{\it Example}: For $A = \mathcal{P}(\mathbb{R})$, we can define the relation $EQ_{\mathbb{R}}$ on $A$ as 
\[
\{ (X_1, X_2 ) \in\mathcal{P}(\mathbb{R})  \times \mathcal{P}(\mathbb{R}) ~|~ |X_1| = |X_2| \}
\]

\vfill

{\it Example}: Let $R_{(\textbf{mod } n)}$ be the set of all pairs of integers $(a, b)$ such that $(a \textbf{ mod } n = b \textbf{ mod } n)$.
Then $a$ is {\bf congruent to} $b$ \textbf{mod} $n$ means $(a, b) \in R_{(\textbf{mod } n)}$. A common notation is to write this as $a \equiv b (\textbf{mod } n)$.


$R_{(\textbf{mod } n)}$ is a relation on the set $\underline{\phantom{\mathbb{Z}}\hspace{20em}}$


Some example elements of $R_{(\textbf{mod } 4)}$ are: 

\vfill

\end{document}